\section{The Binary-Field Polynomial IOP}
\label{sec:piop}

The arithmetization hands us a UAIR whose constraints split into a linear part
(ideal-membership) and a nonlinear part (column Hadamard products). This section
proves the linear part. It is the binary-field specialisation of the
Zinc\texttt{+} polynomial interactive oracle reduction: a commitment, an ideal
check, an evaluation reading $\psia$, and a sumcheck, composing to reduce ``all
ideal-membership constraints hold at all rows'' to a batch of multilinear
evaluation claims at a common point. The Hadamard part is deferred to
\Cref{sec:hadamard}; both meet at the commitment (\Cref{sec:commitment}).

\subsection{Pipeline}
\label{sec:pipeline}

The reduction runs on one Fiat--Shamir transcript through
\[
  \underbrace{\text{commit}}_{\text{Step 0}}
  \to
  \underbrace{\text{ideal check}}_{\text{Step 2}}
  \to
  \underbrace{\psia}_{\text{Step 3}}
  \to
  \underbrace{\text{sumcheck}}_{\text{Step 4}}
  \to
  \underbrace{\text{multipoint eval}}_{\to\, r_0}
  \to
  \underbrace{\text{open}}_{\text{Step 7}},
\]
with the Hadamard discharge of \Cref{sec:hadamard} riding Steps 2--4 --- it is
one more ideal-membership family --- and folding into the multipoint eval.
\emph{Step 1 is empty:} the general Zinc\texttt{+} compiler begins by projecting
$\Q[X]$ or $\Z[X]$ to a prime field $\F_{q_0}[X]$, but the trace here is already
over $\Ftwo[X]\subset\K[X]$, so there is no prime to sample and no projection to
perform. This vacancy is the first concrete saving of the all-$\Ftwo$ choice.

\subsection{Step 0: commitment, and Step 1's vacancy}
\label{sec:commit-step}

The prover commits the primary witness columns with the hash-based $\Ftwo[X]$
commitment of \Cref{sec:commitment}: each column is arranged as a matrix, its
rows encoded by an $\Ftwo$-linear code that does not widen $\cellring$, and the
codewords Merkle-hashed into a single root absorbed into the transcript. Public
columns are not committed (the verifier evaluates their MLEs locally); virtual
columns are not committed (they are reconstructed from sources). For now the
commitment is a black box that will, at the end, certify the evaluation claims
this section produces.

\subsection{Step 2: the ideal check}
\label{sec:ideal-check}

The linear constraints are ideal-membership assertions: for each family $\ell$
with generator $g_\ell\in\{X,\,0\}$ and each row $i$, a fixed
$\Ftwo[X]$-combination $Q_\ell(T_i)$ of cells and row-shifts must satisfy
$Q_\ell(T_i)\in\ideal{g_\ell}$. The ideal check batches \emph{all} of these ---
across families and all $N_{\mathrm{rows}}$ rows --- following the
ideal-batching stage of the Zinc\texttt{+} compiler. It draws a row-evaluation
point $\rho\in\K^{\nu}$ and a batching randomness, and forms, per family, the
combined polynomial
\[
  E_\ell = \sum_{i\in\{0,1\}^{\nu}} \eq_i(\rho)\,Q_\ell(T_i)\;\in\;\K[X],
\]
with the obligation $E_\ell\in\ideal{g_\ell}$. Because $\ideal{g_\ell}$ is a
$\K[X]$-submodule, honest data yields $E_\ell\in\ideal{g_\ell}$; a single
cheating row is caught, with the usual MLE soundness, by the random $\rho$.

\begin{remark}[The $\Ftwo$-native linear ideal check]
\label{rem:native-ic}
Every SHA-256 constraint is degree one in the trace cells, so the ideal check
runs in its linear lane: the $E_\ell$ are $\Ftwo$-linear images of the committed
cells, assembled by evaluating each constraint's bit-MLEs at $\rho$ over $\K$,
with no polynomial-product machinery. The membership obligations specialise to
$E_\ell$ having zero constant term ($g_\ell = X$) or $E_\ell = 0$ outright
($g_\ell = 0$).
\end{remark}

\subsection{Step 3: the evaluation reading $\psia$}
\label{sec:psi-alpha}

The combined claim still lives in $\K[X]$ --- it speaks of the cell variable. The
prover draws $\alpha\in\K$ and applies $\psia$ (substitution $X = \alpha$) to
every cell, producing the projected columns $\mle{\psia(w_g)}$, and evaluates
each $E_\ell$ at $\alpha$. Soundness for the linear part rests on the following
commutation.

\begin{lemma}[Linear constraints survive $\psia$]
\label{lem:linear-survives}
Let $Q$ be a fixed $\Ftwo[X]$-linear combination of cells (and row-shifts) and
$g\in\Ftwo[X]$. For any trace $T$ and any $\alpha\in\K$:
\begin{enumerate}[leftmargin=2em,label=(\roman*)]
  \item $\psia\bigl(Q(T)\bigr) = Q\bigl(\psia(T)\bigr)$, with $\psia$ applied
    cellwise and $Q$ read with its $\Ftwo$ coefficients in $\K$; and
  \item if $Q(T)\in\ideal{g}$ then $\psia(Q(T)) = g(\alpha)\cdot u$ for some
    $u\in\K$; in particular $g(\alpha)=0$ gives $\psia(Q(T))=0$.
\end{enumerate}
\end{lemma}

\begin{proof}
(i) $\psia$ is $\Ftwo$-linear and $Q$ has $\Ftwo$ coefficients, so $\psia$
commutes with the combination; row-shifts are positional and unaffected.
(ii) $Q(T) = g\cdot h$ for some $h$, and $\psia$ is evaluation at $\alpha$, hence
multiplicative: $\psia(Q(T)) = g(\alpha)\,\psia(h)$.
\end{proof}

Part (i) is exact, in \emph{both} directions, and that is what makes the
projected trace a faithful target: \emph{completeness} is immediate (a true
identity among cells projects to a true identity among the projected columns),
and \emph{soundness} is its converse read through randomness --- a \emph{false}
membership survives $\psia$ only if its nonzero degree-$<\wbit$ cell polynomial
vanishes at the random $\alpha$, which happens with probability
$\le(\wbit-1)/|\K|$ (Schwartz--Zippel). Part (ii) is how the obligations
discharge downstream: an LSB constraint ($g = X$) leaves a controlled
$\alpha$-multiple, an equality ($g = 0$) an exact $\K$-equality.

\begin{remark}[Why $\AND$ does \emph{not} survive]
\label{rem:and-not-survive}
\Cref{lem:linear-survives}(i) fails for the Hadamard product: $\psia$ is an
algebra map for the \emph{polynomial} product, but
$\psia(\bp{u}\had\bp{v}) = \sum_b u_b v_b\alpha^b$ is not
$\psia(\bp{u})\,\psia(\bp{v}) = \sum_{b,b'} u_b v_{b'}\alpha^{b+b'}$. A product
of readings is not the reading of a Hadamard product. This is exactly why the
Hadamard constraints are not equalities under $\psia$; \Cref{sec:hadamard} reads
them through a \emph{different} functional $\psiz$ --- at the same point
$\alpha$ --- under which they become an ideal-membership claim with the
bit-product structure exposed rather than scrambled.
\end{remark}

\subsection{Step 4: the sumcheck}
\label{sec:sumcheck}

After the reading the linear obligations are a zerocheck: a batched constraint
MLE must vanish on the hypercube, weighted by $\eq(\cdot;\rho)$. The prover
batches the projected columns by a challenge $\gamma\in\K$ into
$\mle{C}(x) = \sum_g \gamma^{g}\,\mle{\psia(w_g)}(x)$, and the claim
\[
  \sum_{x\in\{0,1\}^{\nu}} \eq(x;\rho)\,\mle{C}(x) = 0
\]
encodes, by \Cref{lem:linear-survives}, that all batched linear constraints
hold. This is a degree-$2$ statement, and a standard sumcheck reduces it in
$\nu$ rounds to a single point $r^*\in\K^{\nu}$ with per-column claims
$a_g = \mle{\psia(w_g)}(r^*)\in\K$. The prover forms each round's message with
$\eq(\cdot;\rho)$ factored out, so the per-slot work is a product against
$\mle{C}$ alone. The Hadamard discharge of \Cref{sec:hadamard} adds one further
(degree-$3$) group to this same sumcheck, reduced to the same $r^*$ and formed
the same way.

\subsection{Multipoint evaluation: a single point $r_0$}
\label{sec:multipoint}

The sumcheck leaves claims at $r^*$, but the AIR reads several previous rows
(the shift-register layout of \Cref{sec:air}), so the verifier also owes
\emph{pointed-shift} claims: evaluations of a source column at a point shifted by
a row offset $\Delta$. And the Hadamard discharge produces its own claims, the
$\psiz$-projections of the operand columns at the same $r^*$. A multipoint
evaluation argument collects all of these and reduces them to one point: given
the claims at $r^*$ and a list of pointed-shift claims, a single sumcheck
produces a common point $r_0\in\K^{\nu}$ at which every involved column --- the
$\psia$-projected trace and, for the Hadamard path, the $\psiz$-projected trace
appended to it --- carries one evaluation claim. The commitment then opens every
column at the single $r_0$.

\subsection{What the PIOP proves, and what is left}
\label{sec:piop-handoff}

At the end of Steps 2--4 and the multipoint eval, the verifier holds the
commitment root and $\alpha$, the common point $r_0$, and, for each committed
column, evaluation claims at $r_0$ for its $\psia$-image (the linear claims) and,
when the Hadamard path runs, its $\psiz$-image. These are exactly the inputs the
commitment's evaluation argument expects. By \Cref{lem:linear-survives},
certifying the $\psia$-claims certifies the entire linear part of the UAIR ---
rotations, message schedule, modular-addition LSBs, and boundaries. What remains
uncertified is the nonlinear part: the $14$ column Hadamard products. The next
section proves those, and arranges its claims to be exactly the $\psiz$-images at
$r_0$ that the same opening will bind --- so the nonlinear part costs no second
opening.
